\metatitle{The 0-Hecke algebra}
\metadescription{The 0-Hecke algebra, its simple and projective modules, and
its relationship with quasisymmetric and noncommutative symmetric functions.}
\metakeywords{0-Hecke algebra,quasisymmetric characteristic,QSym,NSym,
fundamental quasisymmetric functions}

\section[zeroHeckeAlgebra]{The 0-Hecke algebra}

The \defin{0-Hecke algebra}\label{zeroHeckeAlgebraDefinition} \(H_n(0)\) in
type \(A\) is the \(\setC\)-algebra generated by
\(T_i\), \(i=1,\dotsc,n-1\), subject to
\[
T_i^2 = T_i, \qquad
T_iT_{i+1}T_i = T_{i+1}T_iT_{i+1}, \qquad
T_jT_i = T_iT_j \text{ whenever } |i-j|\geq 2.
\]
The braid relations are the same as for the symmetric group, while the
quadratic relation is idempotent rather than involutive.  If
\(w=s_{i_1}\dotsm s_{i_\ell}\) is a
\hyperref[reducedWord]{reduced word}, then
\[
  T_w\coloneqq T_{i_1}\dotsm T_{i_\ell}
\]
is independent of the chosen reduced word.  Thus
\(\{T_w:w\in\symS_n\}\) is a basis, and \(\dim H_n(0)=n!\).

\begin{example}
For \(n=3\), the algebra \(H_3(0)\) has basis
\[
  T_e,\quad T_1,\quad T_2,\quad T_1T_2,\quad T_2T_1,\quad T_1T_2T_1.
\]
The braid relation identifies \(T_1T_2T_1\) with \(T_2T_1T_2\), just as in
the symmetric group.
\end{example}

\subsection[zeroHeckeRepresentationTheory]{Representation theory}

The representation theory of \(H_n(0)\) parallels the representation theory
of \(\symS_n\), but with compositions replacing partitions
\cite{KrobThibon1997}.  The simple \(H_n(0)\)-modules are one-dimensional and
are indexed by compositions \(\alpha\vDash n\), or equivalently by subsets of
\([n-1]\).  We denote them by \(F^\alpha\).

\begin{array}{lll}
\toprule
\text{Object} & \text{Group algebra }\setC[\symS_n] & \text{0-Hecke algebra }H_n(0)\\
\midrule
\text{Indexing set for simples} & \text{partitions of }n & \text{compositions of }n\\
\text{Characteristic target} & \spaceSym & \spaceQSym\\
\text{Simple/module basis} & \schurS_\lambda & \gessel_\alpha\\
\text{Projective dual side} & \spaceSym & \mathrm{NSym}\\
\bottomrule
\end{array}

\svgimg[width=0.78\textwidth]{svg-images/qsym-nsym-zero-hecke.svg}{
The QSym--NSym duality as seen through 0-Hecke modules and projectives.}

In the diagram, the blue node represents QSym, the green node represents
NSym, and the orange node represents the \(0\)-Hecke representation category.

The \defin{quasisymmetric characteristic}
\label{quasisymmetricCharacteristic} is the analogue of the
\hyperref[frobeniusCharacteristic]{Frobenius characteristic}.  It sends the
simple module \(F^\alpha\) to the
\hyperref[gessel]{fundamental quasisymmetric function}
\(\gessel_\alpha\).

\begin{example}
For \(n=3\), the compositions and descent sets are
\[
  (3)\leftrightarrow\varnothing,\qquad
  (1,2)\leftrightarrow\{1\},\qquad
  (2,1)\leftrightarrow\{2\},\qquad
  (1,1,1)\leftrightarrow\{1,2\}.
\]
On the simple module \(F^{(1,2)}\), the generator \(T_1\) acts by \(0\) and
\(T_2\) acts by \(1\).  Its quasisymmetric characteristic is
\(\gessel_{(1,2)}\).  This is the 0-Hecke analogue of sending an irreducible
\(\symS_n\)-module to its Frobenius characteristic.
\end{example}

\subsection[zeroHeckeQSymNSym]{QSym and NSym}

The Grothendieck group of finite-dimensional \(H_n(0)\)-modules, summed over
all \(n\), is identified with \(\spaceQSym\).  The Grothendieck group of
projective \(0\)-Hecke modules is identified with
\hyperref[nonComFunctions]{NSym}.  Under induction and restriction, these
Grothendieck groups recover the dual Hopf-algebra structures on QSym and
NSym \cite{KrobThibon1997}.

This gives a representation-theoretic explanation for the
\hyperref[qsymNSymDuality]{duality between QSym and NSym}.  The fundamental
quasisymmetric basis comes from simple modules, while the dual
noncommutative side comes from projectives.  More precisely, indecomposable
projectives correspond to the \hyperref[nonComRibbonBasis]{noncommutative
ribbon basis}, which is dual to the fundamental quasisymmetric basis.
Induction and restriction of 0-Hecke modules give the product and coproduct
on the corresponding Hopf algebras.

\subsection[zeroHeckeSchurLike]{Schur-like bases}

Many quasisymmetric Schur-like bases have 0-Hecke interpretations.  For
example, \name{Vasu Tewari} and \name{Stephanie van Willigenburg} show that
\hyperref[qSchur]{quasisymmetric Schur functions} arise as quasisymmetric
characteristics of certain \(H_n(0)\)-modules
\cite{TewariWilligenburg2015}.  Similar representation-theoretic models
appear for dual immaculate, extended Schur, and related bases on the
\hyperref[qSchur]{quasisymmetric Schur page}.

Divided-difference and Demazure-type operators also have 0-Hecke flavor.  For
example, the operators used to construct \hyperref[key]{key polynomials} and
\hyperref[macdonaldE]{non-symmetric Macdonald polynomials} satisfy braid and
idempotent-type relations closely related to \(H_n(0)\).

\subsection[zeroHeckeOtherTypes]{Other types}

The type \(B\) 0-Hecke algebra replaces \(\symS_n\) by the
\hyperref[permutationsTypeB]{hyperoctahedral group}.  \name{Young-Hun Kim}
and \name{Dominic Searles} construct type \(B\) poset modules whose
quasisymmetric characteristics give a representation-theoretic interpretation
of type \(B\) \(P\)-partition enumerators \cite{KimSearles2026x}.

This is one instance of a broader pattern: once a Coxeter group has a
\hyperref[coxeterGroups]{Coxeter presentation}, one can form a 0-Hecke
algebra by replacing \(s_i^2=1\) with an idempotent relation.
